\question[5]
Die Kryptologie unterscheidet die vier Sicherheitsziele \textit{Vertraulichkeit}, \textit{Integrit\"at},
\textit{Authentizit\"at} und \textit{Verbindlichkeit}. \\
a. Erl\"autere kurz, was mit den Sicherheitszielen \textit{Integrit\"at} und \textit{Verbindlichkeit} gemeint ist. \\
b. Ordne den folgenden Situationen jeweils das verletzte Sicherheitsziel zu: \\
\hspace*{1em} (1) Eve liest eine unverschl\"usselte E-Mail von Alice an Bob mit. \\
\hspace*{1em} (2) Mallory f\"angt eine \"Uberweisung ab und \"andert den Empf\"anger. \\
\hspace*{1em} (3) Mallory schickt Bob eine Bestellung, die scheinbar von Alice stammt. \\
\hspace*{1em} (4) Alice bestreitet nachtr\"aglich, einen von ihr verschickten Vertrag verfasst zu haben.

\begin{solutionbox}{6cm}
a. \textit{Integrit\"at}: Die empfangene Nachricht wurde von keiner dritten Person unbemerkt ver\"andert
(manipuliert). \\
\textit{Verbindlichkeit}: Der Urheber einer Nachricht kann nachtr\"aglich nicht bestreiten, sie verfasst
zu haben. \\
b. (1) Vertraulichkeit \quad (2) Integrit\"at \quad (3) Authentizit\"at \quad (4) Verbindlichkeit
\end{solutionbox}
